\section{Exact support for quadratic blocks}\label{sec:quadratic}

If every feature $f_j$ is a quadratic polynomial, then every support
problem~\eqref{eq:support} asks for the minimum of a single quadratic
function over the block domain. This section gives two exact rational
algorithms for that problem. The first works for any bounded polytope of
small dimension; the second works for quadratic stars of any size whose
constraints couple the center with one leaf at a time. Both return an
attained minimizer and a rational minimum, so they serve as exact support
oracles and give support inequalities that are tight.

\subsection{Bounded polytopes in fixed dimension}\label{sec:polytope}

Let $P=\{x\in\R^d:Ax\le b\}$ be a bounded polytope given by $m$ rational
rows, which include finite variable bounds, and let
$q(x)=\tfrac12x\T Hx+c\T x+c_0$ with rational data and $H=H\T$. The matrix
$H$ may be indefinite or singular, $P$ may be lower-dimensional, and rows
may be redundant or implied equalities; an equality is given by two
opposite rows. For a set $I$ of rows, $A_I$ and $b_I$ denote the
corresponding submatrix and subvector, and
\begin{equation}\label{eq:kkt-system}
M_I=\begin{pmatrix}H&A_I\T\\ A_I&0\end{pmatrix},\qquad
M_I\begin{pmatrix}x\\ \mu\end{pmatrix}=\begin{pmatrix}-c\\ b_I\end{pmatrix}.
\end{equation}
The \emph{candidates} are the points $x$ obtained from all sets $I$ with
$\abs I\le d$ for which $M_I$ is nonsingular and $Ax\le b$ holds. The
multipliers $\mu$ are not sign-restricted: the system expresses
stationarity on the affine hull of a face, not a KKT condition.

\begin{theorem}[Exact quadratic support on a polytope]\label{thm:polytope}
If $P\ne\emptyset$, some candidate minimizes $q$ over $P$; hence the least
candidate value is $\min_Pq$, and it is attained at a rational point. If
there is no candidate, then $P=\emptyset$. Both statements remain true if
only sets $I$ of linearly independent rows are enumerated.
\end{theorem}

The proof (Appendix~\ref{app:polytope}) chooses a global minimizer in the
relative interior of a face of smallest dimension and shows that $H$ is
positive definite on the face's direction space, so that a nonsingular
system~\eqref{eq:kkt-system} of independent active rows has this minimizer
as its solution. The sign of the multipliers and the inertia of $H$ need not
be tested, because extra candidates, such as saddle points of faces, are
feasible and cannot produce a value below the minimum. Dependent rows never
give a nonsingular $M_I$, so duplicate, redundant and opposite rows are
harmless. Two parts of the argument cannot be dropped. Boundedness is needed:
for $P=\{x\in\R:x\ge0\}$ and $q=-x$ the enumeration returns $0$, and for
$P=\R$ and $q=x$ it returns no candidate. The choice of a minimizer on a face
of minimal dimension is also needed: for $q=(x-y)^2$ on $[0,1]^2$, the
minimizers $(t,t)$ with $0<t<1$ give only singular systems, and the
enumeration finds the minimizers $(0,0)$ and $(1,1)$ instead.

\begin{corollary}[Support of a quadratic feature graph]\label{cor:quad-features}
If $f_1,\dots,f_p$ are quadratic polynomials with rational coefficients and
$(a,b)$ is rational, the support value $h_P(a,b)$ is the rational number
obtained by applying Theorem~\ref{thm:polytope} to $q=a\T x+b\T F$, and the
support inequality is tight at the returned minimizer.
\end{corollary}

There are at most $N(m,d)=\sum_{j=0}^d\binom mj$ sets $I$. For each,
Gaussian elimination on a system of order at most $2d$ and the feasibility
test take $O(d^3+md)$ arithmetic operations, so the enumeration uses
$O(N(m,d)(d^3+md))$ operations. With fraction-free elimination, every
intermediate number is a minor of an integer matrix of order at most $2d$,
and Hadamard's inequality bounds the bit length of every candidate and of
the minimum by $O(d(\langle q\rangle+\langle A\rangle+\log d))$, where
$\langle q\rangle$ is the encoding length of $q$ and $\langle A\rangle$ the
largest encoding length of a row (Appendix~\ref{app:polytope}). For fixed
$d$ the algorithm is therefore strongly polynomial, and the minimum has
polynomial encoding length, in line with the result that quadratic
programming is in NP \citep{Vavasis1990}. The method is total enumeration of
faces \citep[Section~2.9]{Murty1997}; what we add is the treatment of the
degenerate cases that a certificate must cover.

Without a bound on the dimension the problem is hard. The following facts
are folklore \citep[Section~2.3]{BurerLetchford2009}; Appendix~\ref{app:polytope}
gives the short proofs.

\begin{proposition}[Hardness in unrestricted dimension]\label{prop:maxcut}
Computing $\min_{x\in[0,1]^n}q(x)$ for rational quadratics $q$ is strongly
NP-hard, even when $q$ is multilinear with integer coefficients of absolute
value at most $n$. Deciding whether a given inequality $q(x)\ge t$ holds on
$[0,1]^n$ is coNP-complete.
\end{proposition}

Unless NP${}={}$coNP, the validity of an exact support cut for a quadratic
block of unrestricted size therefore has no short certificate that can be
checked in polynomial time, and a certificate that is replayed by repeating
the enumeration, as ours is, is natural for this class. Useful exact oracles
must restrict the dimension, as in Theorem~\ref{thm:polytope}, or the
structure, as in the next subsection. Restricting the interaction graph alone
is not enough: box-constrained quadratic minimization is strongly NP-hard
already at treewidth two, and quartic minimization already on paths
\citep[Theorems~2 and~3]{DelPiaKhajavirad2026}.

\subsection{Quadratic stars with center--leaf constraints}\label{sec:star}

Section~\ref{sec:composition} shows that the blocks worth merging are often
stars: several pairs that share one variable. For stars there is an exact
algorithm whose cost is nearly linear in the size of the block. Let $y$ be
the center and $x_1,\dots,x_k$ the leaves, and consider
\begin{equation}\label{eq:star}
q(y,x)=c_0+c_1y+c_2y^2+\sum_{i=1}^k\bigl(\alpha_ix_i^2+(\beta_iy+\gamma_i)x_i\bigr)
\end{equation}
with rational coefficients of arbitrary signs. The domain $\Omega$ is given by
finite rational bounds on all variables and by rational affine rows, each of
which involves $y$ and \emph{at most one} leaf; an equality is given as two
rows. Every scalarization $a\T(y,x)+b\T F(y,x)$ of a vector $F$ of quadratics
with this interaction pattern has the form~\eqref{eq:star}. Rows that couple
two leaves are excluded, because with them the problem is NP-hard already
without a center: $\min\sum_ix_i(1-x_i)$ over $0\le x\le1$ and
$\sum_ia_ix_i=b$ is zero exactly when the subset-sum instance $(a,b)$ is
feasible \citep{Karp1972}; compare \citet{MoreVavasis1990}.

For fixed $y$, the rows of leaf $i$ restrict $x_i$ to an interval
$[L_i(y),U_i(y)]$, where $L_i$ is the maximum of the affine lower bounds and
$U_i$ the minimum of the affine upper bounds that these rows and the leaf
bounds define. Thus $L_i$ is convex and $U_i$ concave, both piecewise
affine. The set $J_i=\{y:L_i(y)\le U_i(y)\}$ is an interval because $L_i-U_i$
is convex. Given $y$, the leaves are constrained independently, so the
projection of $\Omega$ onto $y$ is the interval $I$ obtained by intersecting
the bounds of $y$, the rows in $y$ alone, and all $J_i$. If $I$ is empty, so
is $\Omega$.

\begin{theorem}[Exact support on constrained stars]\label{thm:star}
Let $m$ be the number of rows that involve a leaf and $m_0$ the number of
rows in $y$ alone, and suppose $\Omega$ is nonempty. Then $\min_\Omega q$ is
attained at a rational point, and it can be computed, together with a
minimizer and a partition of $I$ into at most $2m+4k+1$ intervals on each of
which every leaf has an affine minimizer, with $O(m_0+(m+k)\log(m+k))$
arithmetic operations and comparisons. All numbers involved have bit length
polynomial in the input length; the breakpoints have bit length bounded
independently of $k$.
\end{theorem}

\begin{proof}
Define $\varphi_i(y)=\min\{\alpha_ix_i^2+(\beta_iy+\gamma_i)x_i:L_i(y)\le x_i\le U_i(y)\}$
for $y\in I$, so that $\min_\Omega q=\min_{y\in I}V(y)$ with
$V(y)=c_0+c_1y+c_2y^2+\sum_i\varphi_i(y)$. Let $s_i=m_i+2$ be the number of
affine functions that define $L_i$ and $U_i$, where $m_i$ is the number of
rows of leaf $i$; together $L_i$ and $U_i$ have at most $s_i-2$ breakpoints.

If $\alpha_i>0$, the minimizer is $\xi_i=\max\{L_i,\min\{v_i,U_i\}\}$ with
the affine function $v_i(y)=-(\beta_iy+\gamma_i)/(2\alpha_i)$. The function
$v_i-L_i$ is concave and $v_i-U_i$ convex, so each has at most two isolated
zeros, and $\xi_i$ is affine between consecutive points of a set of at most
$s_i+2$ breakpoints. If $\alpha_i\le0$, a minimizer is an endpoint of the
interval, and
\begin{equation}\label{eq:endpoint-factor}
\psi_i(U_i)-\psi_i(L_i)=(U_i-L_i)\bigl(\alpha_i(U_i+L_i)+\beta_iy+\gamma_i\bigr),
\qquad\psi_i(t)=\alpha_it^2+(\beta_iy+\gamma_i)t .
\end{equation}
The first factor is nonnegative on $I$. On each of the at most $s_i-1$
intervals where $L_i$ and $U_i$ are affine, the second factor is affine and
has at most one zero, which decides the preferred endpoint; this gives at
most $2s_i-3$ breakpoints. Taking all breakpoints together splits $I$ into at
most $\sum_i2s_i+1=2m+4k+1$ intervals on which every leaf rule is affine and
$V(y)=q(y,\xi(y))$ is a quadratic polynomial in $y$.

Each $\varphi_i$ is continuous on $I$, since it is the minimum of a continuous
function over an interval whose endpoints depend continuously on $y$. Hence
$V$ agrees with the quadratic of each interval on its closure, and the
minimum of $V$ is the least value of these quadratics at the interval
endpoints and, where the leading coefficient is positive, at interior
stationary points. All breakpoints and all endpoints of $I$ are zeros of
affine functions with rational coefficients, so every candidate is rational,
and the minimizing $y$ together with $x_i=\xi_i(y)$ is a rational minimizer.

For the operation count, intersecting the bounds of $y$ with the $m_0$ rows
in $y$ alone costs $O(m_0)$. The upper envelope of $s_i$ lines is computed by
sorting slopes and one stack-based scan in $O(s_i\log s_i)$ operations, and
$J_i$, the breakpoints and the rules of leaf $i$ in $O(s_i)$ more. Sorting all
breakpoints costs $O((m+k)\log(m+k))$. A sweep over the sorted breakpoints
maintains the three coefficients of $V$; at each breakpoint only the leaves
whose rule changes are updated, so the sweep takes $O(m+k)$ operations. Bit
lengths are bounded in Appendix~\ref{app:star}.
\end{proof}

The bound on the number of intervals is tight up to a constant factor: a
single concave leaf whose bounds zigzag forces $2s_i-3$ rule changes, and $k$
hinge leaves force $k$ more, so the partition can need $\Theta(m+k)$
intervals. On instances with $m=\Theta(k)$, the operation bound is optimal up
to a constant factor for algebraic computation trees: by Ben-Or's theorem,
applied to a family of instances with $3r$ leaves and $r$ rows whose decision
set has $r!$ connected components, deciding whether the minimum is at least a
given value requires $\Omega((m+k)\log(m+k))$ operations
(Appendix~\ref{app:star}).

\paragraph{Relation to forest algorithms.}
\citet{DelPiaKhajavirad2026} solve box-constrained quadratic programs on
forests with $O(n^2)$ arithmetic operations. After an affine scaling of each
variable to $[0,1]$, a box-constrained star is such a forest, and their
dynamic program rooted at the center specializes to the box-constrained case
of Theorem~\ref{thm:star}; in that case our contribution is only the
$O(k\log k)$ count, which replaces their $O(k^2)$ merge at the root. Rows
that couple the center with a leaf are outside their model. A shear
$x_i\mapsto x_i-\rho y$ turns a leaf domain that is a parallelogram with two
vertical sides back into a box and keeps the star structure, but triangles,
trapezoids with nonparallel sides and polygons with more sides cannot be
reduced in this way. Such rows also break the structural property on which
their algorithm rests: the conditional value of a leaf is then no longer a
minimum of affine functions of $y$ over a fixed set, and it can have convex
kinks. For example, for $x\in[0,1]$ with rows $x\ge y-\tfrac12$ and
$x\ge\tfrac12-y$, the minimum of $x$ is $\abs{y-\tfrac12}$. The SOC
descriptions of star hulls by \citet{DeyKhajavirad2025} assume sign patterns
of the square terms and box domains, and the LP approximations of
\citet{BienstockMunoz2018} allow constraints but are approximate. Our
algorithm conditions on the center, as nonserial dynamic programming does
\citep{BerteleBrioschi1972}, and solves a one-parameter quadratic problem for
each leaf, a special case of multiparametric quadratic programming
\citep{Bemporad2002}; \citet{BhathenaEtAl2025} use parametric dynamic
programs of this kind on trees for convex quadratic problems with
indicators. To our knowledge, no exact algorithm for nonconvex stars with
center--leaf rows has been published.

\begin{remark}[Deeper trees]\label{rem:trees}
The argument uses that every vertex other than the center is a leaf. On the
path $x_1$--$x_2$--$x_3$ over $[0,1]^3$ with
$f=(x_2-\tfrac18)x_1+x_2^2-x_2+x_2x_3+x_3^2-x_3$, the conditional value of the
subtree $\{x_1,x_2\}$ given $x_3=t$ is
$\min\{-\tfrac18,\,-\tfrac14(1-t)^2\}$, which switches at the irrational point
$t=1-\sqrt2/2$, although $\min f=-\tfrac38$ is rational
(Appendix~\ref{app:star}). Rooted at $x_2$, the same path is a star, and
Theorem~\ref{thm:star} gives a rational partition; the obstruction concerns
nesting conditional values along a tree of depth two.
\citeauthor{DelPiaKhajavirad2026} handle such breakpoints for boxes with an
explicit algebraic representation; we do not know whether edge-local affine
rows can be added to deeper trees with a polynomial algorithm.
\end{remark}

The certificate of a star support value records the partition of $I$, the
affine rule of every leaf on every piece, the quadratic of $V$ on every piece
and the candidates; a checker rebuilds the partition from the rows and
coefficients and compares. The star oracle in our implementation is simpler
than the algorithm of the proof: it splits the center interval at all
pairwise intersections of each leaf's bound lines and recomputes every leaf
rule on every piece. This costs $O((m+k)^3)$ operations, stores $\Theta(k)$
rules per piece, and can produce more pieces than Theorem~\ref{thm:star}
requires. It certified the star cuts of our first implementation, whose
stars had at most four leaves (Appendix~\ref{app:campaigns12}); in the
experiments of Section~\ref{sec:computations} every block had at most four
variables, and the polytope oracle certified every quadratic cut.
Section~\ref{sec:star-bench} compares the two star algorithms on random
stars with up to $10^4$ leaves.
